% Irodov, section 2.1 "Equation of the Gas State. Processes",
% original numbered problems 1--11 (printed pp.62--63; PDF pp.71--72).
% Section 2 uses separate ORIGINAL numbering; do not concatenate to 1..388.

\thermoproblem{1}{62}
一容器体积为 $V=30\,\mathrm L$，
其中装有 $0^\circ\mathrm C$ 的理想气体。
在温度保持不变的情况下放出一部分气体，
容器压强降低
$\Delta p=0.78\,\mathrm{atm}$。
气体在标准状况下的密度为
$\rho_0=1.3\,\mathrm{g/L}$。
求放出的气体质量。
\exsolution
设标准状况压强
$p_0=1.0\,\mathrm{atm}$、
温度 $T_0=273.15\,\mathrm K$。
题中容器温度同为标准温度，
因此气体的密度与压强成正比，
由理想气体状态方程，放出质量对应
\[
\Delta m=\frac{M}{RT_0}\Delta p\,V.
\]
利用标准状况
$\rho_0=Mp_0/(RT_0)$，
消去未知摩尔质量：
\[
\boxed{\Delta m=\rho_0 V\frac{\Delta p}{p_0}
=(1.3\,\mathrm{g/L})(30\,\mathrm L)(0.78)
\simeq30\,\mathrm g}.
\]
\excheck 质量与容器体积以及压强降幅成正比；
不需要单独知道气体的摩尔质量。

\thermoproblem{2}{62}
两只体积相同的容器由单向阀相连；
当两侧压强差达到或超过
$\Delta p=1.10\,\mathrm{atm}$ 时，
阀门允许气体从高压侧流向低压侧。
最初一只容器为真空，
另一只装有
$t_1=27^\circ\mathrm C$、
$p_1=1.00\,\mathrm{atm}$
的理想气体。
随后两只容器升温到
$t_2=107^\circ\mathrm C$。
求最初真空容器达到的压强。
\exsolution
两只容器总气体摩尔数不变且最终同温。
设每只容积为 $V$，
初始有气容器的气体满足
$nR T_1=p_1V$。
升温后若气体仍全部在原容器，其压强为
$p_\Sigma=p_1 T_2/T_1$。
阀门达到停止流动的临界状态时，
两侧压强之差为 $\Delta p$。
令最初真空侧压强为 $p$，
原有气体一侧压强为 $p+\Delta p$。
因为两容器体积相等、末温相同，
\[
(p+\Delta p)V+pV=nRT_2
=p_1V\frac{T_2}{T_1}.
\]
解得
\[
\boxed{p=\frac12\left(
p_1\frac{T_2}{T_1}-\Delta p\right)}.
\]
取 $T_1\simeq300\,\mathrm K$、
$T_2\simeq380\,\mathrm K$，
得到
\[
\boxed{p\simeq0.083\,\mathrm{atm}}.
\]
\excheck 原书简答记为约
$0.10\,\mathrm{atm}$；严格代入本版
$p_1=1.00\,\mathrm{atm}$、$\Delta p=1.10\,\mathrm{atm}$
给出约 $0.083\,\mathrm{atm}$。
这个差异可理解为答案过度粗略取整，
已在勘误中说明。

\thermoproblem{3}{62}
体积为 $V=20\,\mathrm L$ 的容器内
装有氢气与氦气的理想混合气，
温度 $t=20^\circ\mathrm C$、
压强 $p=2.0\,\mathrm{atm}$，
总质量 $m=5.0\,\mathrm g$。
求氢气质量与氦气质量之比。
\exsolution
令氢气、氦气质量分别为 $m_H,m_{\rm He}$，
其摩尔质量近似为
$M_H=2\,\mathrm{g/mol}$、
$M_{\rm He}=4\,\mathrm{g/mol}$。
由气体分压相加与理想气体状态方程：
\[
\frac{pV}{RT}=\frac{m_H}{M_H}
+\frac{m_{\rm He}}{M_{\rm He}},
\qquad
m_H+m_{\rm He}=m.
\]
令总摩尔数
$n=pV/(RT)$，以克为单位，
$n=m_H/2+(5-m_H)/4=(m_H+5)/4$，
所以
\[
m_H=4n-5,\quad m_{\rm He}=10-4n.
\]
取
$R=0.082057\,\mathrm{L\,atm/(mol\,K)}$、
$T=293.15\,\mathrm K$，
$n\simeq1.663\,\mathrm{mol}$，
\[
\boxed{\frac{m_H}{m_{\rm He}}
=\frac{4n-5}{10-4n}\simeq0.50}.
\]
\excheck 两组分的\emph{摩尔数}并非质量比例；
需用各自摩尔质量换算。

\thermoproblem{4}{62}
容器中混有
$m_1=7.0\,\mathrm g$ 的氮气与
$m_2=11\,\mathrm g$ 的二氧化碳，
温度 $T=290\,\mathrm K$，
总压强 $p_0=1.0\,\mathrm{atm}$。
把两气体视为理想气体，
求混合气体的质量密度 $\rho$。
\exsolution
混合气体总摩尔数等于各组分摩尔数之和：
\[
n=\frac{m_1}{M_1}+\frac{m_2}{M_2}.
\]
由理想气体方程
$p_0V=nRT$，
得到容积
$V=(m_1/M_1+m_2/M_2)RT/p_0$。
因此密度为
\[
\boxed{\rho=\frac{m_1+m_2}{V}
=\frac{p_0(m_1+m_2)}
{RT(m_1/M_1+m_2/M_2)}}.
\]
取 $M_1=28\,\mathrm{g/mol}$、
$M_2=44\,\mathrm{g/mol}$，
则两组分各为 $0.25$ mol，
总质量 $18$ g，总摩尔数 $0.50$ mol，
\[
\boxed{\rho\simeq1.5\,\mathrm{g/L}}.
\]
\excheck 混合气平均摩尔质量为
$M_{\rm av}=18/0.50=36\,\mathrm{g/mol}$，
用 $\rho=p_0M_{\rm av}/(RT)$ 可复核。

\thermoproblem{5}{63}
容器体积 $V=7.5\,\mathrm L$，
其中装有温度 $T=300\,\mathrm K$
的理想气体混合物：
氧气 $\nu_1=0.10$ mol、
氮气 $\nu_2=0.20$ mol、
二氧化碳 $\nu_3=0.30$ mol。
求：(a) 混合气总压强；
(b) 用于方程 $pV=(m/M)RT$
的平均摩尔质量 $M$。
\exsolution
气体总摩尔数
$\nu=\nu_1+\nu_2+\nu_3=0.60$ mol。
道尔顿分压定律及理想气体方程给
\[
\boxed{p=\frac{\nu RT}{V}
=\frac{0.60(0.082057)(300)}{7.5}
\simeq1.97\,\mathrm{atm}\approx2.0\,\mathrm{atm}}.
\]
混合物总质量
$m=\nu_1M_1+\nu_2M_2+\nu_3M_3$，
平均摩尔质量
\[
\boxed{M=\frac{m}{\nu}
=\frac{0.10(32)+0.20(28)+0.30(44)}
{0.60}
\simeq36.7\,\mathrm{g/mol}}.
\]
\excheck 质量加权与\emph{摩尔数加权}不同；
用组分摩尔分数计算平均摩尔质量。

\thermoproblem{6}{63}
一竖直密闭圆筒由无摩擦轻活塞分成
上下两腔，各装 1 mol 空气。
在温度 $T_0=300\,\mathrm K$ 的平衡态，
上腔体积是下腔体积的
$\eta=4.0$ 倍。
当两腔达到新的共同温度时，
体积比变为 $\eta'=3.0$，
求新的温度 $T$。
\exsolution
设圆筒总容积为 $V_\Sigma$、
活塞横截面积为 $S$、
活塞质量为 $m$。
当上、下腔体积比为 $\eta$，
\[
V_u=\frac{\eta V_\Sigma}{1+\eta},
\qquad V_d=\frac{V_\Sigma}{1+\eta}.
\]
轻活塞竖直力平衡
$(p_d-p_u)S=mg$，
理想气体状态方程
$p_d=RT/V_d$、
$p_u=RT/V_u$，
因此
\[
\frac{mg}{S}
=\frac{RT}{V_\Sigma}
\frac{\eta^2-1}{\eta}.
\]
右侧的 $mg/S,V_\Sigma$ 不变，
故
\[
\frac{T_0(\eta^2-1)}{\eta}
=\frac{T(\eta'^2-1)}{\eta'}.
\]
解出
\[
\boxed{T=T_0
\frac{\eta'(\eta^2-1)}
{\eta(\eta'^2-1)}}
=\boxed{300\cdot\frac{3\cdot15}{4\cdot8}
\simeq422\,\mathrm K}.
\]
\excheck 温度升高后两侧压强差仍由活塞重量决定，
体积比趋向更接近1；题目没有必要单独知道活塞质量。

\thermoproblem{7}{63}
体积为 $V$ 的密闭容器通过活塞式抽气泵抽气，
每次活塞行程在吸气时形成有效体积 $\Delta V$。
假设等温理想气体，
求把容器压强降低 $\eta$ 倍
所需抽气行程数 $n$。
\exsolution
一次抽气时，容器与泵室连通后
气体等温占据体积 $V+\Delta V$，
而总摩尔数暂时不变。
根据玻意耳定律，
若抽气前容器压强为 $p_j$，
抽气后重新隔离容器时
\[
p_{j+1}=p_j\frac{V}{V+\Delta V}.
\]
重复 $n$ 次：
\[
p_n=p_0\left(\frac{V}{V+\Delta V}\right)^n.
\]
要求 $p_n=p_0/\eta$，
取对数得到
\[
\boxed{n=\frac{\ln\eta}
{\ln(1+\Delta V/V)}}.
\]
\excheck 实际抽气次数须向上取整；
此式假定泵无死区容积、漏气、阀门回流或温度变化。

\thermoproblem{8}{63}
体积为 $V$ 的容器最初压强为 $p_0$，
由抽速为常数 $C$ 的真空泵等温抽气。
抽速定义为\emph{当前腔室压强下}
单位时间抽走的气体体积。
求容器压强与时间 $t$ 的关系。
\exsolution
以气体摩尔数 $n=pV/(RT)$ 描述容器状态。
单位时间抽走的摩尔数为
$\mathrm d n_{\rm out}/\mathrm dt=Cp/(RT)$。
故
\[
\frac{\mathrm d}{\mathrm dt}\left(\frac{pV}{RT}\right)
=-\frac{Cp}{RT}.
\]
因为 $T,V$ 不变，
\[
\frac{\mathrm dp}{\mathrm dt}
=-\frac CV\,p.
\]
分离变量并用初态 $p(0)=p_0$，
得
\[
\boxed{p(t)=p_0\exp(-Ct/V)}.
\]
\excheck 这里抽速 $C$ 并不是恒定\emph{抽出摩尔数}；
随着压强降低，单位时间抽出的摩尔数减少，
所以压强指数下降。

\thermoproblem{9}{63}
容积为 $V=87\,\mathrm L$ 的真空室，
抽气泵在当前压强下的抽速恒为
$C=10\,\mathrm{L/s}$。
求真空室压强降低 $\eta=1000$ 倍
所需时间。
\exsolution
根据上一题等温恒定抽速规律
$p(t)=p_0e^{-Ct/V}$。
令 $p(t)=p_0/\eta$，得
\[
e^{-Ct/V}=\eta^{-1}
\quad\Longrightarrow\quad
\boxed{t=\frac VC\ln\eta}.
\]
代入数值：
\[
\boxed{t=\frac{87\,\mathrm L}{10\,\mathrm{L/s}}
\ln1000
\simeq60\,\mathrm s\simeq1.0\,\mathrm{min}}.
\]
\excheck $V/C$ 具有时间单位；
未计入泵的极限真空度、出气和漏气。

\thermoproblem{10}{63}
原书图2.1中，一根竖直光滑管上下段横截面积不同，
上、下各有一只活塞，并用不可伸长细线相连。
两活塞之间密封 1 mol 理想气体。
上活塞面积比下活塞大
$\Delta S=10\,\mathrm{cm^2}$，
两活塞总质量 $m=5.0\,\mathrm{kg}$，
外界大气压 $p_0=1.0\,\mathrm{atm}$。
求要使两活塞整体向上平移
$\ell=5.0\,\mathrm{cm}$，气体应升温多少。
\exsolution
对两个活塞作为整体列竖直静力平衡。
活塞间气压 $p$ 对两活塞的净向上推力为
$p(S_u-S_d)=p\Delta S$，
外界大气压产生净向下力 $p_0\Delta S$，
总重力为 $mg$，故
\[
(p-p_0)\Delta S=mg
\quad\Longrightarrow\quad
p=p_0+\frac{mg}{\Delta S}.
\]
在活塞随绳共同上移 $\ell$ 时，
气体体积增加
$\Delta V=(S_u-S_d)\ell=\Delta S\,\ell$。
准静态过程中 $p$ 不变。
由一摩尔理想气体状态方程，
\[
R\Delta T=p\Delta V
=(p_0\Delta S+mg)\ell,
\]
因此
\[
\boxed{\Delta T=\frac{(p_0\Delta S+mg)\ell}{R}}.
\]
用 $\Delta S=10^{-3}\,\mathrm{m^2}$、
$p_0\simeq1.013\times10^5\,\mathrm{Pa}$、
$m=5.0\,\mathrm{kg}$、
$\ell=0.050\,\mathrm m$：
\[
\boxed{\Delta T\simeq0.90\,\mathrm K}.
\]
\excheck \textbf{题面图像复核：}
原书 PDF 纯文本抽取曾把
$\ell=5.0\,\mathrm{cm}$ 错读为 $50\,\mathrm{cm}$。
回看原版放大的题面图像，小数点清楚存在；
按 $\ell=0.050\,\mathrm m$ 计算约为
$0.90\,\mathrm K$，与本版书末简答一致。
该例\emph{不是原书答案错误}。

\thermoproblem{11}{63}
一摩尔理想气体发生下列准静态过程，
求沿该过程能够达到的最高温度，
其中 $p_0,\alpha,\beta$ 是正常数：
(a) $p=p_0-\alpha V^2$；
(b) $p=p_0e^{-\beta V}$。
\exsolution
一摩尔理想气体满足
$T=pV/R$。
(a) 温度函数
$T(V)=(p_0V-\alpha V^3)/R$，
对 $V$ 求导：
\[
\frac{\mathrm dT}{\mathrm dV}
=\frac{p_0-3\alpha V^2}{R}.
\]
内部极大值在
$V_*=\sqrt{p_0/(3\alpha)}$，
其处二阶导数
$-6\alpha V_*/R<0$。
用
$\alpha V_*^2=p_0/3$：
\[
\boxed{T_{\max}^{(a)}
=\frac{2p_0}{3R}
\sqrt{\frac{p_0}{3\alpha}}}.
\]
(b) $T(V)=p_0Ve^{-\beta V}/R$，
求导
\[
\frac{\mathrm dT}{\mathrm dV}
=\frac{p_0}{R}e^{-\beta V}(1-\beta V).
\]
唯一极大值在 $V_*=1/\beta$，
故
\[
\boxed{T_{\max}^{(b)}
=\frac{p_0}{e\beta R}}.
\]
\excheck 极值必须位于物理允许的
$p>0,V>0$ 区域：
(a) 的 $V_*$ 满足 $p=2p_0/3>0$；
(b) 所有 $V>0$ 时 $p>0$。
